\section{Experimental Evaluation}\label{sec:experiments}

% S = Simulation setup
\subsection{Dataset and Setup}
We use the daily average prices of three vegetables---cabbage, bok choy, and
sponge gourd---in the 13 public retail markets of Taichung
City~\cite{taichung2026retail}, from 2025-10-04 to 2026-10-02, in NT\$ per
catty (600\,g). The publisher uses a price of zero to mean that no price was
recorded; such entries are treated as missing, never as a $-100\%$ change.
Each sample is a (market, vegetable, day $t$) triple whose target is the
relative change $r_h$ of \eqref{eq:delta}, paired on the exact date $t+h$.

The split is chronological and grouped by date, so all markets of one day fall
in the same part: training from 2025-10-04 to 2026-05-08, validation from
2026-05-10 to 2026-07-21, and a held-out test period from 2026-07-22 to
2026-10-02 (73 dates). Assuming a one-day publication delay, a label matures at
$t+h+1$; 1{,}386 rows whose labels crossed a split boundary (695 and 691 at the
two boundaries) were purged to prevent leakage~\cite{kaufman2012leakage}. After
requiring complete features and matured labels, the test sets contain
$n=2{,}085$ samples for $h=7$ and $n=1{,}859$ samples for $h=14$.

Each sample has 12 numeric features built only from data published by
$t-1$---the last visible price and its age, 7-, 14-, and 28-day moving
averages, 14-day volatility and range, 7- and 14-day momentum, weekday, and
month---plus a one-hot encoding of the vegetable; the market identity is not
used. Models are implemented in Python~3.12 with
scikit-learn~1.9.1~\cite{pedregosa2011scikit}.

% Q = Quality and quantity (metrics), R = Related methods to compare
\subsection{Baselines and Metrics}
Three simple baselines are used, following the forecasting practice of always
comparing with naive methods~\cite{makridakis2020m4}:
\begin{itemize}
\item \textit{No-change (persistence):} $\hat{r}_h=0$, i.e., the price stays
as it is.
\item \textit{Momentum:} the latest visible change of the market price over
the past $h$ days continues (missing values treated as $0$).
\item \textit{Train mean by crop:} the mean training-set $r_h$ of each
vegetable and horizon.
\end{itemize}

\textit{Direction macro-F1.} Each change is mapped to a class
$c\in\mathcal{C}=\{\text{down},\text{flat},\text{up}\}$ with the flat band
$\tau=5\%$ of Section~\ref{sec:method}. The band was frozen before the test
period was examined: prices are integers in NT\$ per catty, typically 25--60,
so a one-dollar step is already 2--4\% and a change within $\pm5\%$ does not
alter a purchase decision. With $F_1^{(c)}$ the F1-score of class $c$,
\begin{equation}
\text{macro-F1}=\frac{1}{|\mathcal{C}|}\sum_{c\in\mathcal{C}}F_1^{(c)} .
\label{eq:f1}
\end{equation}
Macro averaging prevents the dominant ``flat'' class from hiding errors on the
rare but important up and down moves.

\textit{Relative-change MAE.}
\begin{equation}
\text{MAE}=\frac{1}{n}\sum_{i=1}^{n}\bigl|r_i-\hat{r}_i\bigr| .
\label{eq:mae}
\end{equation}
Higher macro-F1 and lower MAE are better; MAE is measured on the market's
relative change, not in NT\$ and not on store prices.

% T = Tuning parameters
\subsection{Parameter Settings}
Ridge, RandomForest, and HistGradientBoosting regressors and two direction
classifiers were trained, and the model with the best validation macro-F1 was
selected for each horizon: a RandomForest with 300 trees and at least 5 samples
per leaf for $h=7$, and a HistGradientBoosting regressor with 200 iterations and
learning rate 0.05 for $h=14$, both with random seed 42. Each model predicts
$\hat{r}_h$, which is then classified with $\tau$. The mode-B weight is fixed at
$\alpha=0.3$. The test period was used only once and played no part in tuning.

% U = Useful results
\subsection{Results}
Table~\ref{tab:comparison} compares the selected models with the baselines.
Both models carry some directional signal: their macro-F1 (0.422 at 7 days and
0.417 at 14 days) is well above that of the no-change baseline (0.309 and
0.295), which predicts ``flat'' for every sample. However, neither beats the
momentum baseline (0.423 and 0.431), and their relative-change MAE (0.066 and
0.098) is worse than that of the no-change baseline (0.044 and 0.075). Recall on
the rare up class is only 10.1\% (16 of 158) at 7 days and 8.6\% (23 of 268) at
14 days. We report these negative results as they are.

% TODO(author): copy the verified values from the DIC3 experiment logs and
% delete every \tbd. Do NOT reuse numbers from the AI-generated infographic:
% they contradict the written finding (see paper/README.md).
\begin{table}[t]
\caption{Results on the held-out test sets (values pending verification)}
\label{tab:comparison}
\centering
\footnotesize
\begin{tabular}{@{}llcc@{}}
\toprule
Horizon & Method & Macro-F1 $\uparrow$ & MAE $\downarrow$ \\
\midrule
7 days ($n=2{,}085$) & No-change & \tbd & \tbd \\
 & Momentum & \tbd & \tbd \\
 & RF, mode A & \tbd & \tbd \\
 & RF, mode B & \tbd & \tbd \\
\midrule
14 days ($n=1{,}859$) & No-change & \tbd & \tbd \\
 & Momentum & \tbd & \tbd \\
 & HGB, mode A & \tbd & \tbd \\
 & HGB, mode B & \tbd & \tbd \\
\bottomrule
\end{tabular}
\end{table}


% V = Verify effectiveness (discussion)
\subsection{Discussion}
\textit{Why are the baselines hard to beat?} In our data, 86.7--92.8\% of
seven-day price pairs have exactly the same price at both ends, depending on
the vegetable, and only 10.5\% of seven-day changes leave the $\pm5\%$ band.
Such sticky series~\cite{nakamura2008five,bils2004some} make ``no change''
close to optimal for absolute error, while momentum captures the few
persistent moves, for example after a typhoon. This agrees with large
forecasting studies in which simple methods remain
competitive~\cite{makridakis2020m4,zeng2023transformers}.

\textit{RQ1.} At both 7 and 14 days, the evidence does not support using the
learned forecasts alone; the momentum signal should be shown to the user next
to them.

\textit{RQ2.} Mode A preserves the predicted direction by construction
(Section~\ref{sec:method}), so it cannot improve the direction metrics on
public data. Whether mode B adds value cannot be tested here, because the
public data contain no store prices; this requires matured same-store prices
collected after each query, which has not yet started.

\textit{RQ3.} The web prototype and the MCP tools answer a query with dated
7- and 14-day results, and every number shown comes from a tool call. This
shows that low-effort querying is technically feasible; a user study is still
needed to measure usefulness.

\textit{Threats to validity.} The target is the public market price, which is
only a proxy for the price in a specific store; the data cover three
vegetables in one city over one year; and the direction metric depends on the
5\% band.
